\problemname{Bitland Coin}
The central bank of Bitland just launched its official cryptocurrency Bitland Coin (BLC). To mine a block of BLC
you need to find a sufficiently small value of $\texttt{hash}(x \,.\, k)$ where $\,.\,$ is bit-concatenation, $x$ is a $28$-bit number you can't change,
$k$ is a $36$-bit number you can choose freely and $\texttt{hash}$ is the function defined below.
A different way of writing the same thing would be $\texttt{hash}(2^{36} x + k)$.

Bitland where confident they had the resources and the expertise to create a hash function better than SHA256.
One of their genious ideas was to use a $64$-bit input and output to save storage space (bits are not as abundant in Bitland as you might expect).
The result of their extensive research, the hash function used in BLC, is defined by the following python code.

\begin{verbatim}
def hash(x):
    # int to list of bits
    def tolist(x):
        return [(x>>i)&1 for i in range(63, -1, -1)]

    # list of bits to int
    def toint(x):
        return sum([x[63-i]<<i for i in range(0, 64)])

    assert 0 <= x < 2**64
    x = tolist(x)
    for i in range(512):
        y = 4*x[i%64] + 2*x[(i+2)%64] + x[(i+10)%64]
        x[i%64] = [1,1,0,0,1,1,0,1][y]
    x = toint(x)
    for i in range(32):
        x += toint(tolist(x)[::-1])
        x ^= 3**40
        if(x >= 2**64): x >>= 1
    assert 0 <= x < 2**64
    return x
\end{verbatim}

You want to get rich by mining this coin. BLC automatically changes its difficulty (how small the hash needs to be). Therefore, you want to write a program which finds a value of $k$ 
to make the hash as small as possible, the smaller the better.

\section*{Input}
Input consists of $10$ lines. The $i$th line contains one integer $x_i$. All $x_i$ are sampled uniformly at random from $[0, 2^{28}-1]$.

\section*{Output}
Output $10$ integers $k_i$, one per line.

\section*{Scoring}
Your program will only be run on one testcase (containing $10$ lines).
Your score for the $i$th line is $10 - \frac{10}{64} \cdot \log_2(\texttt{hash}(2^{36} \cdot x_i + k_i))$. Your total score is the sum of scores
for each line, rounded to two decimal places. The sample testcase is not worth any points.

\section*{Explanation of sample}
The first line of input and output results in the hash $\texttt{hash}(2^{36} \cdot 111131462 + 57474941445) = \texttt{hash}(7636895975021609477) = 112582045000626944$.
The score for the first line is $10 - \frac{10}{64} \cdot \log_2(112582045000626944) \approx 1.14941$. The same procedure for each of the following lines gives a
total score of TODO.
